\originalchapter{电磁场的波动与辐射}\label{ch:em}
\originalsection{场本身的振动}
看见吉他弦, 是反射光把有关弦的信息带到眼睛; 听见吉他, 是空气压强扰动带到耳朵. 光在真空中仍能传播, 所以它不需要像空气那样的物质介质. 变化的是空间每一点的电场 $\boldsymbol E$ 与磁场 $\boldsymbol B$. 本节只建立后续光学所需的波与能流; 完整场论及边界积分推导见本仓库电动力学册第7章“麦克斯韦方程与电磁势”、第8章“能量、动量与角动量”及第9章“电磁波、界面与波导”.

散度 $\nabla\cdot\boldsymbol A$ 描述单位体积的净向外通量, 旋度 $\nabla\times\boldsymbol A$ 描述单位小面面积的环流. 这些局部量与积分量由散度定理和斯托克斯定理连接. 无电荷、无电流的真空区域中,
\[
\nabla\cdot\boldsymbol E=\nabla\cdot\boldsymbol B=0,\quad
\nabla\times\boldsymbol E=-\boldsymbol B_t,\quad
\nabla\times\boldsymbol B=\mu_0\epsilon_0\boldsymbol E_t.
\]
对法拉第式取旋度, 用 $\nabla\times(\nabla\times\boldsymbol E)=\nabla(\nabla\cdot\boldsymbol E)-\nabla^2\boldsymbol E$, 得
\[
\boldsymbol E_{tt}=c^2\nabla^2\boldsymbol E,\quad
\boldsymbol B_{tt}=c^2\nabla^2\boldsymbol B,\qquad c=1/\sqrt{\mu_0\epsilon_0}.
\]
电场的每个分量满足波方程, 但三个分量不能任意独立选择, 还必须满足全部麦克斯韦方程.

设复平面波 $\boldsymbol E=\Rea[\boldsymbol E_0\e^{\ii(\boldsymbol k\cdot\boldsymbol r-\omega t)}]$. 散度条件给 $\boldsymbol k\cdot\boldsymbol E_0=0$, 说明电场横向; 旋度条件给 $\boldsymbol B_0=\boldsymbol k\times\boldsymbol E_0/\omega$. 真空色散关系为 $\omega=c|\boldsymbol k|$, 故 $\boldsymbol B=\hat{\boldsymbol k}\times\boldsymbol E/c$. 场彼此垂直且同相, 传播方向为 $\boldsymbol E\times\boldsymbol B$. 常量静场是另一个可能存在的解, 并不属于这对传播波.

单位体积场能与能流密度是 $u=\epsilon_0E^2/2+B^2/(2\mu_0)$ 与 $\boldsymbol S=\boldsymbol E\times\boldsymbol B/\mu_0$. 对单一真空平面波, $u=\epsilon_0E^2$, $\boldsymbol S=cu\hat{\boldsymbol k}$, 平均强度为 $I=\epsilon_0c|\boldsymbol E_0|^2/2$. 对左右行波形成的驻波, $E,B$ 并不同相, 局部能量在电磁两种形式间交换, 不能用单一行波的 $B=E/c$ 比较每一点每一时刻的幅度.

两束等幅反向线偏振波取 $E_y^{\pm}=E_0\cos(kx\mp\omega t)$, 相应 $B_z^+=E_y^+/c$, $B_z^-=-E_y^-/c$. 相加得到
\[
E_y=2E_0\cos kx\cos\omega t,\qquad
B_z=\frac{2E_0}{c}\sin kx\sin\omega t.
\]
所以 $u=2\epsilon_0E_0^2(\cos^2kx\cos^2\omega t+\sin^2kx\sin^2\omega t)$, $S_x=\epsilon_0cE_0^2\sin(2kx)\sin(2\omega t)$. 时间平均为 $\langle u\rangle=\epsilon_0E_0^2$, $\langle S_x\rangle=0$. 电场腹对应磁场节点; 在 $t=0$ 全部场能是电能, 在 $t=\pi/(2\omega)$ 全部是磁能. 直接求导可检查 $u_t+\partial_xS_x=0$, 瞬时能流把能量在相邻位置间转移.

\originalsection{介质中的有效响应}
介质中的电荷受束缚, 外场使原子中的正负电荷产生微小相对位移. 单位体积电偶极矩是极化强度 $\boldsymbol P$, 磁偶极矩是磁化强度 $\boldsymbol M$. 极化随时间变化对应束缚电流 $\boldsymbol J_P=\boldsymbol P_t$, 磁化还给 $\boldsymbol J_M=\nabla\times\boldsymbol M$. 将 $\boldsymbol J_{\rm bound}=\boldsymbol P_t+\nabla\times\boldsymbol M$ 代入微观安培--麦克斯韦方程并移项, 得 $\nabla\times\boldsymbol H=\boldsymbol J_{\rm free}+\boldsymbol D_t$. 因而介质极化影响传播, 并非凭空改变真空常数. 定义 $\boldsymbol D=\epsilon_0\boldsymbol E+\boldsymbol P$, $\boldsymbol H=\boldsymbol B/\mu_0-\boldsymbol M$. 对均匀、各向同性、线性且在考察频带近似无损的介质, 使用 $\boldsymbol D=\epsilon\boldsymbol E$, $\boldsymbol B=\mu\boldsymbol H$, 波速 $v=1/\sqrt{\mu\epsilon}$, 折射率 $n=c/v$. 磁性很弱时 $\mu\simeq\mu_0$, $n^2\simeq\epsilon/\epsilon_0$.

物质响应通常随频率变化, 因而 $n=n(\omega)$, $k=n(\omega)\omega/c$. 有耗散时 $\epsilon$ 及 $k$ 可为复数, 表示相位传播同时幅度衰减; 不能在强吸收区仍以实 $v$ 描述所有性质. 对无损非色散介质, 平均能流为 $I=|E_0|^2/(2Z_{\rm em})$, 其中 $Z_{\rm em}=\sqrt{\mu/\epsilon}$. 色散介质的储能需要考虑频率响应, 不把简单的 $\epsilon E^2/2$ 当作一般公式.

\originalsection{加速电荷与远场}
匀速运动的孤立电荷不向无限远持续辐射, 而受加速度的电荷会发出传播扰动. 对非相对论速度、局域电荷运动且观察距离 $r$ 远大于源的尺寸, 电场中随 $1/r$ 下降的辐射项为
\begin{equation}\label{eq:radiation}
\boldsymbol E_{\rm rad}(\boldsymbol r,t)
=-\frac{q}{4\pi\epsilon_0c^2r}\boldsymbol a_\perp(t-r/c),\quad
\boldsymbol a_\perp=\boldsymbol a-(\boldsymbol a\cdot\hat{\boldsymbol r})\hat{\boldsymbol r},\quad
\boldsymbol B_{\rm rad}=\frac1c\hat{\boldsymbol r}\times\boldsymbol E_{\rm rad}.
\end{equation}
延迟时间 $t-r/c$ 是光从源到观察点的传播时间. 此式是电动力学册第10章“偶极辐射与运动电荷”的远场结果; 它的使用条件在这里保留, 不将形象的场线弯折图冒充完整场论证明. 辐射场按 $1/r$ 下降, 强度按 $1/r^2$ 下降, 球面总功率因此保持有限. 近场项不贡献无穷远的辐射功率. 若要在有限距离把此辐射项近似为总场, 对典型角频率 $\omega$ 还需处于 $r\gg c/\omega$ 的辐射区; 单独的静电场没有传播波的坡印廷通量.

以 $\boldsymbol a$ 为极轴, 夹角 $\theta$ 的每单位立体角辐射功率为 $q^2a^2\sin^2\theta/(16\pi^2\epsilon_0c^3)$. 由 $\int\sin^2\theta\dd\Omega=2\pi\int_0^\pi\sin^3\theta\dd\theta=8\pi/3$, 得拉莫尔功率 $P=q^2a^2/(6\pi\epsilon_0c^3)$. 电荷沿 $z$ 轴振动 $z=d\cos\omega t$ 时, 轴上 $a_\perp=0$ 没有这一阶偶极辐射; 垂直于轴的方向最强且线偏振. 若电荷在 $xy$ 平面作椭圆运动 $x=\sqrt2A\cos\omega t$, $y=A\sin\omega t$, 沿 $x$ 或 $y$ 观察只有一个横分量, 线偏振; 沿 $z$ 观察两个横分量均存在且相差四分之一周期, 椭圆偏振. 光的偏振于是直接联系到源的运动与观察方向.
